% =====================================================================
%  PRÉAMBULE COMMUN — Carnet d'entraînement, Fiche 2 (Nomenclature)
%  (inséré physiquement dans chaque document pour qu'il soit autonome)
% =====================================================================
\documentclass[11pt,a4paper]{article}

% ---------- Encodage & langue ----------
\usepackage[utf8]{inputenc}
\usepackage[T1]{fontenc}
\usepackage{lmodern}
\usepackage[french]{babel}

% ---------- Mathématiques & chimie ----------
\usepackage{amsmath,amssymb,mathtools}
\usepackage{icomma}
\usepackage[version=4]{mhchem}
\usepackage{siunitx}
\sisetup{output-decimal-marker={,},per-mode=symbol}

% ---------- Graphismes & mise en page ----------
\usepackage{xcolor}
\usepackage{colortbl}
\usepackage{tikz}
\usepackage[most]{tcolorbox}
\usepackage{enumitem}
\usepackage{array}
\usepackage{booktabs}
\usepackage{multicol}
\usepackage{fancyhdr}
\usepackage{titlesec}
\usepackage[hidelinks]{hyperref}
\usetikzlibrary{arrows.meta,positioning,calc,shapes.geometric,shapes.misc,
  fit,backgrounds,decorations.pathreplacing}
\usepackage[a4paper,margin=2.2cm,headheight=28pt,headsep=10pt]{geometry}

% =====================================================================
%  PALETTE (charte « Chimie » — mauve/prune du cours)
% =====================================================================
\definecolor{primary}{RGB}{146,95,161}
\definecolor{primarydark}{RGB}{92,55,108}
\definecolor{defcol}{RGB}{0,114,178}
\definecolor{thmcol}{RGB}{0,140,120}
\definecolor{formulacol}{RGB}{198,93,9}
\definecolor{remarkcol}{RGB}{120,94,200}
\definecolor{attncol}{RGB}{200,40,55}
\definecolor{excol}{RGB}{0,135,90}
\definecolor{methcol}{RGB}{90,90,100}
\definecolor{metalcol}{RGB}{198,150,40}
\definecolor{nonmetcol}{RGB}{0,140,150}
\definecolor{oxycol}{RGB}{210,55,45}
\definecolor{hydcol}{RGB}{60,120,210}
\definecolor{catcol}{RGB}{200,70,120}
\definecolor{ancol}{RGB}{40,110,190}
\definecolor{saltcol}{RGB}{120,90,160}

% =====================================================================
%  STYLE COMMUN DES BOÎTES
% =====================================================================
\tcbset{
  boxbase/.style={enhanced,breakable,boxrule=0.9pt,arc=2.5pt,
     left=3mm,right=3mm,top=2.3mm,bottom=2.3mm,
     fonttitle=\bfseries,coltitle=white,
     toptitle=0.6mm,bottomtitle=0.6mm},
}

\newtcolorbox{definition}[1][]{%
  boxbase,colback=defcol!5,colframe=defcol,
  title={\textsc{Définition}\ifx\relax#1\relax\relax\else\textmd{ --- \itshape #1}\fi}}
\newtcolorbox{regle}[1][]{%
  boxbase,colback=thmcol!6,colframe=thmcol,
  title={\textsc{Règle}\ifx\relax#1\relax\relax\else\textmd{ --- \itshape #1}\fi}}
\newtcolorbox{formule}[1][]{%
  boxbase,colback=formulacol!6,colframe=formulacol,
  title={\textsc{Méthode-clé}\ifx\relax#1\relax\relax\else\textmd{ --- \itshape #1}\fi}}
\newtcolorbox{remarque}[1][]{%
  boxbase,colback=remarkcol!6,colframe=remarkcol,
  title={\textsc{Remarque}\ifx\relax#1\relax\relax\else\textmd{ : \itshape #1}\fi}}
\newtcolorbox{attention}[1][]{%
  boxbase,colback=attncol!6,colframe=attncol,
  title={\textsc{Attention}\ifx\relax#1\relax\relax\else\textmd{ : \itshape #1}\fi}}
\newtcolorbox{exemple}[1][]{%
  boxbase,colback=excol!5,colframe=excol,
  title={\textsc{Exemple}\ifx\relax#1\relax\relax\else\textmd{ : \itshape #1}\fi}}
\newtcolorbox{methode}[1][]{%
  boxbase,colback=methcol!7,colframe=methcol,sharp corners,
  title={\textsc{Méthode}\ifx\relax#1\relax\relax\else\textmd{ : \itshape #1}\fi}}
\newtcolorbox{astuce}[1][]{%
  boxbase,colback=primary!6,colframe=primary,
  title={\textsc{Astuce concours}\ifx\relax#1\relax\relax\else\textmd{ : \itshape #1}\fi}}

% ---- Exercice (numéroté) ----
\newtcolorbox[auto counter]{exercice}[1][]{%
  boxbase,colback=primary!4,colframe=primarydark,
  title={\textsc{Exercice}~\thetcbcounter\ifx\relax#1\relax\relax\else\textmd{ --- \itshape #1}\fi}}

% ---- Corrigé (numéroté) ----
\newtcolorbox[auto counter]{corrige}[1][]{%
  boxbase,colback=excol!4,colframe=excol,
  title={\textsc{Corrigé}~\thetcbcounter\ifx\relax#1\relax\relax\else\textmd{ --- \itshape #1}\fi}}

% =====================================================================
%  EN-TÊTES ET PIEDS DE PAGE
% =====================================================================
\pagestyle{fancy}
\fancyhf{}
\fancyhead[L]{\small\textcolor{primarydark}{\textbf{Fiche 2} — Nomenclature}}
\fancyhead[R]{\small\textcolor{primarydark}{\textsc{Carnet d'entraînement}}}
\fancyfoot[C]{\small\thepage}
\renewcommand{\headrulewidth}{0.8pt}
\renewcommand{\headrule}{\hbox to\headwidth{\color{primary}\leaders\hrule height \headrulewidth\hfill}}

\titleformat{\section}{\Large\bfseries\color{primarydark}}{\thesection}{0.6em}{}
\titleformat{\subsection}{\large\bfseries\color{primary}}{\thesubsection}{0.6em}{}
\titleformat{\subsubsection}{\normalsize\bfseries\color{primary!80!black}}{}{0em}{}

% ---- Bandeau de titre réutilisable : \bandeau{Thème}{Sous-titre} ----
\newcommand{\bandeau}[2]{%
  \begin{center}
  \begin{tikzpicture}
  \node[rectangle,rounded corners=7pt,fill=primary,text=white,
        minimum width=16.0cm,minimum height=1.7cm,align=center,inner sep=8pt]
    {{\Large\bfseries #1}\\[3pt]{\normalsize #2}};
  \end{tikzpicture}
  \end{center}
  \vspace{0.3cm}}
\begin{document}
\bandeau{Thème 2 — Moyen}{Du nom à la formule : croisement des valences --- Corrigé}

\begin{corrige}[écrire un sel neutre à partir de son nom]
\begin{enumerate}[label=\alph*),leftmargin=1.8em,topsep=1pt,itemsep=1pt]
  \item sulfure de sodium : \ce{S^2-}$+$\ce{Na+} $\to$ \ce{Na2S}. \quad $2(+1)+(-2)=0$. \checkmark
  \item nitrate de calcium : \ce{NO3-}$+$\ce{Ca^2+} $\to$ \ce{Ca(NO3)2}. \quad $(+2)+2(-1)=0$. \checkmark
  \item phosphate d'aluminium : \ce{PO4^3-}$+$\ce{Al^3+} $\to$ \ce{AlPO4} (valences égales à 3 $\Rightarrow$
        indices 1). \quad $(+3)+(-3)=0$. \checkmark
  \item sulfate de potassium : \ce{SO4^2-}$+$\ce{K+} $\to$ \ce{K2SO4}. \quad $2(+1)+(-2)=0$. \checkmark
  \item carbonate de sodium : \ce{CO3^2-}$+$\ce{Na+} $\to$ \ce{Na2CO3}. \quad $2(+1)+(-2)=0$. \checkmark
\end{enumerate}
\end{corrige}

\begin{corrige}[quand la valence du métal est précisée]
\begin{enumerate}[label=\alph*),leftmargin=1.8em,topsep=1pt,itemsep=1pt]
  \item nitrate de fer(III) : \ce{Fe^3+}$+$\ce{NO3-} $\to$ \ce{Fe(NO3)3}.
  \item sulfate de cuivre(II) : \ce{Cu^2+}$+$\ce{SO4^2-} $\to$ \ce{Cu2(SO4)2} $\Rightarrow$ simplifié \ce{CuSO4}.
  \item chlorure d'étain(IV) : \ce{Sn^4+}$+$\ce{Cl-} $\to$ \ce{SnCl4}.
  \item carbonate de fer(II) : \ce{Fe^2+}$+$\ce{CO3^2-} $\to$ \ce{Fe2(CO3)2} $\Rightarrow$ simplifié \ce{FeCO3}.
  \item phosphate de fer(III) : \ce{Fe^3+}$+$\ce{PO4^3-} $\to$ \ce{FePO4} (3 et 3 $\Rightarrow$ indices 1).
\end{enumerate}
\end{corrige}

\begin{corrige}[le piège des suffixes -ite / -ate]
\begin{enumerate}[label=\alph*),leftmargin=1.8em,topsep=1pt,itemsep=1pt]
  \item nitr\emph{ite} \ce{NO2-} $\to$ \ce{NaNO2} \quad|\quad nitr\emph{ate} \ce{NO3-} $\to$ \ce{NaNO3}
        (un oxygène de plus).
  \item sulf\emph{ite} \ce{SO3^2-} $\to$ \ce{K2SO3} \quad|\quad sulf\emph{ate} \ce{SO4^2-} $\to$ \ce{K2SO4}
        (un oxygène de plus).
  \item hypochlor\emph{ite} \ce{ClO-} $\to$ \ce{NaClO} \quad|\quad perchlor\emph{ate} \ce{ClO4-} $\to$ \ce{NaClO4}
        (trois oxygènes de plus).
\end{enumerate}
Les deux anions de chaque paire ont la \emph{même} valence, donc seul le nombre d'oxygènes change ; le
croisement, lui, reste identique.
\end{corrige}

\begin{corrige}[cation \emph{et} anion polyatomiques : le thiosulfate d'ammonium]
Cation : ammonium \ce{NH4+} (valence 1). Anion : thiosulfate \ce{S2O3^2-} (valence 2).
Croisement : la valence 2 de l'anion descend sur le cation, la valence 1 du cation reste sur l'anion :
\[ \ce{NH4+} + \ce{S2O3^2-} \;\longrightarrow\; \ce{(NH4)2S2O3}. \]
Parenthèses autour de \ce{NH4} (polyatomique, $\times 2$). Vérif : $2(+1)+(-2)=0$. \checkmark \quad
$\boxed{\ce{(NH4)2S2O3}}$. \\
À distinguer du \emph{sulfite} d'ammonium \ce{(NH4)2SO3} : le thiosulfate \ce{S2O3^2-} possède un atome de
soufre \emph{de plus} (à la place d'un oxygène) que le sulfite \ce{SO3^2-}.
\end{corrige}

\begin{corrige}[traduire un énoncé d'annales en formules]
\begin{itemize}[leftmargin=1.5em,topsep=1pt,itemsep=1pt]
  \item phosphate de sodium : \ce{PO4^3-}$+$\ce{Na+} $\to$ \ce{Na3PO4}.
  \item bromure de strontium : \ce{Br-}$+$\ce{Sr^2+} $\to$ \ce{SrBr2}.
  \item phosphate de strontium : \ce{PO4^3-}$+$\ce{Sr^2+} $\to$ \ce{Sr3(PO4)2} (croisement 2 et 3).
  \item bromure de sodium : \ce{Br-}$+$\ce{Na+} $\to$ \ce{NaBr}.
\end{itemize}
\emph{(Pour information : l'équation équilibrée serait $\ce{2 Na3PO4 + 3 SrBr2 -> Sr3(PO4)2 + 6 NaBr}$,
mais l'exercice ne demandait que les formules.)}
\end{corrige}
\end{document}
